\documentclass[11pt]{amsart}

\usepackage[T1]{fontenc}
\usepackage{lmodern}
\usepackage{microtype}
\usepackage{amsmath,amssymb,amsthm}
\usepackage[hidelinks]{hyperref}

\hypersetup{
  pdftitle={Counterexamples to the Wide-Interval Case of Conjecture 4 of Gautschi and Leopardi}
}

\newtheorem{theorem}{Theorem}
\newtheorem{corollary}[theorem]{Corollary}
\theoremstyle{remark}
\newtheorem{remark}[theorem]{Remark}

\title[Counterexamples to Conjecture 4]
{Counterexamples to the Wide-Interval Case of Conjecture~4 of Gautschi and Leopardi}

\date{}
\subjclass[2020]{Primary 33C45; Secondary 26D05}
\keywords{Jacobi polynomials, monotonicity, counterexample}

\begin{document}

\begin{abstract}
Gautschi and Leopardi conjectured that, on $0<\theta<\pi$,
validity at $n=1$ of a normalized comparison between consecutive
Jacobi polynomials implies validity for every $n\ge1$. Let
$\beta_*$ be the unique root in $(1,2)$ of
$\beta^3+27\beta^2-10\beta-24=0$; numerically,
$\beta_*=1.1186397525\ldots$. We show that every parameter pair
$(\alpha,\beta)=(0,\beta)$ with $\beta>\beta_*$ is a counterexample:
the comparison holds at $n=1$ for every $\theta\in(0,\pi)$ but fails
at $n=2$ near $\theta=\pi$. For $\beta=2$ we give an exact interior
witness.
\end{abstract}

\maketitle

\section{The counterexamples}

For $\alpha,\beta>-1$, set
\[
\widetilde P_n^{(\alpha,\beta)}(x)
=
\frac{P_n^{(\alpha,\beta)}(x)}{P_n^{(\alpha,\beta)}(1)}.
\]
Gautschi and Leopardi considered
\begin{equation}\label{eq:comparison}
\widetilde P_n^{(\alpha,\beta)}
 \!\left(\cos\frac{\theta}{n}\right)
<
\widetilde P_{n+1}^{(\alpha,\beta)}
 \!\left(\cos\frac{\theta}{n+1}\right).
\end{equation}
Let $\Theta_1^{(\alpha,\beta)}\in(0,\pi)$ be defined by
$P_1^{(\alpha,\beta)}(\cos\Theta_1^{(\alpha,\beta)})=0$. For each
of the intervals $0<\theta<\Theta_1^{(\alpha,\beta)}$ and
$0<\theta<\pi$, Conjecture~4 of Gautschi and Leopardi asserts that
validity of \eqref{eq:comparison} for $n=1$ throughout the interval
implies validity for every $n\ge1$ throughout that interval
\cite[Conjecture~4]{GautschiLeopardi}.

Gautschi and Leopardi report numerical evidence for both branches. In
the passage below, which follows the statement of Conjecture~4, their
(3.6) is \eqref{eq:comparison} and their (3.7) is the pair of intervals
above \cite[Section~3.2]{GautschiLeopardi}:
\begin{quote}
The Matlab script of Section 2.2 is easily adapted to deal with the
conjecture (3.6) for general Jacobi polynomials. When run with the
same data as used to verify Conjecture 3, with $n=100$ and $N=1000$,
similar results were obtained as in Section 3.1, i.e., a strong
indication that (3.6) holds on either interval (3.7) for all
$n\ge1$ whenever $\alpha>1$ and $\beta>-1$, while for $\alpha$ in the
interval $(-1,1)$ the same is true for $\beta$ above a certain curve
that extends from the point $(\alpha,\beta)=(-1,0)$ down to the point
$(\alpha,\beta)=(1,-1)$. For $\beta$ below that curve, the conjecture
fails consistently when $n=1$.
\end{quote}
The sweep used to verify Conjecture~3 takes
$\alpha=0.9,0.8,\dots,-0.9$ and, for each of these, $\beta=20,19,
\dots,0$ and then negative values down to $-0.999$
\cite[Section~3.1]{GautschiLeopardi}; thus $(\alpha,\beta)=(0,2)$ is
one of its nodes. The curve described in the passage runs from
$(\alpha,\beta)=(-1,0)$ down to $(1,-1)$ and so nowhere exceeds
$\beta=0$, whence that node lies well inside the region for which a positive
indication is reported for both intervals.

At that node the failure is not marginal. In the notation of the proof
of Theorem~\ref{thm:family} with $\beta=2$ one has
$\Delta_2(\pi)=-3/16$, and $\Delta_2$ changes sign at
$\theta\approx2.8725$, that is at $\theta\approx0.9144\pi$. On the
grid $\theta_\nu=\nu\pi/1001$, $\nu=1,\dots,1000$, of their routine
for $0<\theta<\pi$ with $N=1000$, all $1000$ nodes satisfy
\eqref{eq:comparison} at $n=1$, the smallest margin there being
$3.7\times10^{-6}$, while $85$ of them violate it at
$n=2$, the first of these being $\nu=916$:
\[
\Delta_2(\theta_{916})=-0.0013,\quad
\Delta_2(\theta_{966})=-0.1032,\quad
\Delta_2(\theta_{1000})=-0.1849.
\]
A plausible explanation is
visible in the verification core printed in their Section~2.2, where
the upper limit of the $\theta$-sweep is set by
\texttt{th1=\allowbreak acos(\allowbreak-1/\allowbreak(2*a+1));} while
the alternative
\texttt{\% th1=pi;} is present but commented out. The script actually
used for Conjecture~4 is described there only as an adaptation of that
one and is not printed, so this explanation is circumstantial rather
than documented.

Koumandos proved \eqref{eq:comparison} for every $n\ge1$ and every
$\theta\in(0,\pi)$ at
$(\alpha,\beta)=(1/2,1/2),(1/2,-1/2),(-1/2,1/2)$ \cite{Koumandos};
that is precisely the branch refuted below, so it is true at those
three pairs and false at $(\alpha,\beta)=(0,\beta)$ for every
$\beta>\beta_*$ (Theorem~\ref{thm:family}). Recent work explicitly
distinguishes this polynomial-value problem from the corresponding
questions for Jacobi zeros
\cite[Section~1]{CastilloRafaeli}. For the zeros the analogous
implication is already known to fail. Extend the notation above by
letting $\cos\Theta_n^{(\alpha,\beta)}$ be the largest zero of
$P_n^{(\alpha,\beta)}$. Conjecture~3 of \cite{GautschiLeopardi}
asserts in the same way that validity at $n=1$ of
$n\Theta_n^{(\alpha,\beta)}<(n+1)\Theta_{n+1}^{(\alpha,\beta)}$
implies validity for every $n\ge1$. Gautschi disproved it on a very
narrow lens-shaped region inside the square
$\{-1<\alpha<-1/2,\ -1/2<\beta<0\}$ and there formulated a revised
parameter domain \cite{Gautschi2008}; the conjectured domains for the
zeros were then extended and modified in
\cite{Gautschi2009a,Gautschi2009b}, with a further remark in
\cite{Gautschi2011}. All four of those papers concern the zeros; none
revisits the polynomial-value implication \eqref{eq:comparison}. To
our knowledge, no counterexample to that implication has previously
been recorded.

\begin{theorem}\label{thm:family}
Let $\beta_*\in(1,2)$ be the unique root in that interval of
\[
q(\beta)=\beta^3+27\beta^2-10\beta-24.
\]
For every $\beta>\beta_*$, the pair $(\alpha,\beta)=(0,\beta)$ is a
counterexample to the $0<\theta<\pi$ branch of Conjecture~4. More
precisely,
\[
P_1^{(0,\beta)}(\cos\theta)
<
P_2^{(0,\beta)}\!\left(\cos\frac{\theta}{2}\right)
\qquad (0<\theta<\pi),
\]
whereas there exists $\varepsilon_\beta>0$ such that
\[
P_2^{(0,\beta)}\!\left(\cos\frac{\theta}{2}\right)
>
P_3^{(0,\beta)}\!\left(\cos\frac{\theta}{3}\right)
\qquad (\pi-\varepsilon_\beta<\theta<\pi).
\]
\end{theorem}

\begin{proof}
Since $P_n^{(0,\beta)}(1)=1$, the normalization is immaterial. From
\[
P_n^{(0,\beta)}(1-2z)
={}_2F_1(-n,n+\beta+1;1;z)
\]
we obtain
\begin{align*}
P_1^{(0,\beta)}(x)
&=1-(\beta+2)z,\\
P_2^{(0,\beta)}(x)
&=1-2(\beta+3)z
  +\frac{(\beta+3)(\beta+4)}{2}z^2,\\
P_3^{(0,\beta)}(x)
&=1-3(\beta+4)z
  +\frac{3(\beta+4)(\beta+5)}{2}z^2
  -\frac{(\beta+4)(\beta+5)(\beta+6)}{6}z^3,
\end{align*}
where $z=(1-x)/2$.

Let $c=\cos(\theta/2)$ and $u=1-c$. Then $0<u<1$, and direct
substitution gives
\begin{equation}\label{eq:first-step}
P_2^{(0,\beta)}(c)-P_1^{(0,\beta)}(\cos\theta)
=
u\left(\beta+1+\frac{\beta^2-\beta-4}{8}u\right).
\end{equation}
The affine factor in parentheses has endpoint values
\[
\beta+1
\quad\text{and}\quad
\frac{\beta^2+7\beta+4}{8},
\]
both positive because $\beta>\beta_*>1$. Thus
\eqref{eq:first-step} is positive for every $0<\theta<\pi$.

Now define
\[
\Delta_2(\theta)
=
P_3^{(0,\beta)}\!\left(\cos\frac{\theta}{3}\right)
-
P_2^{(0,\beta)}\!\left(\cos\frac{\theta}{2}\right).
\]
At the endpoint,
\begin{equation}\label{eq:endpoint}
\Delta_2(\pi)
=
P_3^{(0,\beta)}\!\left(\frac12\right)-P_2^{(0,\beta)}(0)
=-\frac{q(\beta)}{384}.
\end{equation}
Since $q(1)=-6$, $q(2)=72$, and
$q'(\beta)=3\beta^2+54\beta-10>0$ for $\beta\ge1$, the number
$\beta_*$ is well defined and $q(\beta)>0$ for $\beta>\beta_*$. Hence
\eqref{eq:endpoint} is negative. Continuity therefore gives
$\Delta_2(\theta)<0$ whenever $\theta<\pi$ is sufficiently close to
$\pi$.
\end{proof}

\begin{corollary}\label{cor:witness}
At $(\alpha,\beta)=(0,2)$, let
\[
\theta_0=6\arccos\frac78.
\]
Then $0<\theta_0<\pi$, the comparison \eqref{eq:comparison} holds at
$n=1$ for every $\theta\in(0,\pi)$, and it fails at $n=2$ and
$\theta=\theta_0$.
\end{corollary}

\begin{proof}
Put $\varphi=\arccos(7/8)$. Since $7/8>\sqrt3/2$, one has
$0<\varphi<\pi/6$ and hence $0<\theta_0<\pi$. Moreover,
\[
\cos\frac{\theta_0}{2}=\cos(3\varphi)=\frac7{128},
\qquad
\cos\frac{\theta_0}{3}=\cos(2\varphi)=\frac{17}{32}.
\]
Using the formulas above,
\[
P_3^{(0,2)}\!\left(\frac{17}{32}\right)
-P_2^{(0,2)}\!\left(\frac7{128}\right)
=-\frac{6785}{65536}<0.
\]
The $n=1$ assertion follows from Theorem~\ref{thm:family}.
\end{proof}

\begin{remark}[Scope]
The theorem refutes only the $0<\theta<\pi$ branch of
Conjecture~4. It does not decide the
$0<\theta<\Theta_1^{(\alpha,\beta)}$ branch, Conjecture~5 of
\cite{GautschiLeopardi}, or any separate conjecture concerning Jacobi
zeros. Nor is the family $\alpha=0$ minimal: with
\[
\widetilde P_n^{(\alpha,\beta)}(1-2z)
={}_2F_1(-n,n+\alpha+\beta+1;\alpha+1;z),
\]
the same endpoint computation gives
\[
\widetilde P_3^{(\alpha,\beta)}\!\left(\tfrac12\right)
-\widetilde P_2^{(\alpha,\beta)}(0)
=-\tfrac{49}{320},\quad
-\tfrac{167}{37440},\quad
-\tfrac{151}{6720},\quad
-\tfrac{5}{96}
\]
at $(\alpha,\beta)=(-\tfrac12,1)$, $(\tfrac14,\tfrac32)$,
$(\tfrac12,2)$, $(1,3)$ respectively, so the pairs of
Theorem~\ref{thm:family} form a slice of a
two-dimensional region on which \eqref{eq:comparison} fails at $n=2$
near $\theta=\pi$. We work with $\alpha=0$ because there the
normalization is trivial and the witnesses are exactly rational.
\end{remark}

\begin{thebibliography}{9}

\bibitem{GautschiLeopardi}
W.~Gautschi and P.~Leopardi,
\emph{Conjectured inequalities for Jacobi polynomials and their
largest zeros},
Numer. Algorithms \textbf{45} (2007), no.~1--4, 217--230.
\href{https://doi.org/10.1007/s11075-007-9067-5}
{doi:10.1007/s11075-007-9067-5}.

\bibitem{Gautschi2008}
W.~Gautschi,
\emph{On a conjectured inequality for the largest zero of Jacobi
polynomials},
Numer. Algorithms \textbf{49} (2008), 195--198.
\href{https://doi.org/10.1007/s11075-008-9207-6}
{doi:10.1007/s11075-008-9207-6}.

\bibitem{Gautschi2009a}
W.~Gautschi,
\emph{On conjectured inequalities for zeros of Jacobi polynomials},
Numer. Algorithms \textbf{50} (2009), 93--96.
\href{https://doi.org/10.1007/s11075-008-9217-4}
{doi:10.1007/s11075-008-9217-4}.

\bibitem{Gautschi2009b}
W.~Gautschi,
\emph{New conjectured inequalities for zeros of Jacobi polynomials},
Numer. Algorithms \textbf{50} (2009), 293--296.
\href{https://doi.org/10.1007/s11075-008-9230-7}
{doi:10.1007/s11075-008-9230-7}.

\bibitem{Gautschi2011}
W.~Gautschi,
\emph{Remark on ``New conjectured inequalities for zeros of Jacobi
polynomials''},
Numer. Algorithms \textbf{57} (2011), 511.
\href{https://doi.org/10.1007/s11075-010-9442-5}
{doi:10.1007/s11075-010-9442-5}.

\bibitem{Koumandos}
S.~Koumandos,
\emph{On a conjectured inequality of Gautschi and Leopardi for Jacobi
polynomials},
Numer. Algorithms \textbf{44} (2007), no.~3, 249--253.
\href{https://doi.org/10.1007/s11075-007-9098-y}
{doi:10.1007/s11075-007-9098-y}.

\bibitem{CastilloRafaeli}
K.~Castillo and F.~R.~Rafaeli,
\emph{Affine scaling of Jacobi zeros: sharp orderings beyond
Gautschi's conjectures},
arXiv:2608.08258 [math.CA], 2026.

\end{thebibliography}

\end{document}